\section{Integral Weyl comparison from an explicit averaging operator}
\label{sec:integral-weyl}

\status{The following proves the integral Fock comparison used in CL.5--CL.8, and supplies the quantitative isotypic-projector lemma needed by the constrained Gram argument. The background is finite-dimensional rational representation theory of split $\mathrm{GL}_k$ in characteristic zero: complete reducibility, highest weights and the elementary PBW weight bound. The integral claims below are proved separately.}

\subsection{The denominator, and why the operator projects}

Let $\mathcal O$ be a DVR of mixed characteristic $(0,p)$ with fraction field $K$.
Suppose operators $E_{ab}$ satisfy the $\mathfrak{gl}_k$ commutator relations, and let $h\ge0$ be an integer. Write
\[
 D_p(k,h)=\binom{k}{2}\lfloor\log_p(h+k-1)\rfloor,
 \qquad E_{ab}^{[j]}=E_{ab}^j/j!;
 \tag{IW.1}
\]
the cases $k=1$ or $h=0$ below permit exponent zero.
For $z\notin S$, $s=|S|$, define
\[
 \Omega_0=1,\quad \Omega_{q+1}=\sum_{i\in S}E_{zi}\Omega_qE_{iz},\quad
 P_{S,z}=\sum_{q=0}^{h}\frac{(-1)^qs!}{(q+s)!q!}\Omega_q.
\]
Compose these steps while adjoining the $k$ indices successively; call the result $I_h$.
If $M$ is an $\mathcal O$-module stable under every off-diagonal divided power, then
\[
 p^{D_p(k,h)}I_hM\subseteq M. \tag{IW.2}
\]
Indeed, commuting left and right multiplications expands the $q$th summand into divided-power words with coefficients
$(-1)^qs!\prod_i\alpha_i!/(q+s)!$, where $\sum_i\alpha_i=q$.
Its denominator is a product of $s$ binomial coefficients with upper entries at most $h+s$.
Legendre's formula gives $v_p\binom nt\le\lfloor\log_p n\rfloor$.
Thus a step costs at most $s\lfloor\log_p(h+k-1)\rfloor$; sum $s=0,\ldots,k-1$. \textsf{[IW1]}

Suppose all diagonal eigenvalues on the relevant module lie in $[-h,h]$.
If $E_{ab}u=0$ for $a,b\in S$ and $E_{zz}u=0$, the commutator relations give
\[
 E_{jz}\Omega_{q+1}u=\Omega_{q+1}E_{jz}u
 +(q+1)(q+s+1)\Omega_qE_{jz}u\quad(j\in S).
\]
The displayed coefficients telescope to $c_h\Omega_hE_{jz}u$, where $c_h=(-1)^hs!/((h+s)!h!)$.
Each right-hand string inside this expression reaches $z$-weight $-h-1$, so it vanishes.
The opposite root also kills $P_{S,z}u$: for a vector $y$ of both diagonal weights zero with $E_{jz}y=0$,
$E_{jz}E_{zj}^{m+1}y=-m(m+1)E_{zj}^my$; nilpotence and downward induction give $E_{zj}y=0$.
Induction over $S$ proves that $I_h$ sends weight zero to invariants.
Every correction lies in a root-operator image, and the construction commutes with intertwiners.
Consequently on $\operatorname{End}(U)$, for irreducible $U$, it sends a weight-zero $T$ to $(\operatorname{tr}T/\dim U)1$. \textsf{[IW2]}

\subsection{The actual polynomial lattice and its normalized pairing}

Use a square matrix $X=(X_{ij})$ and a partition $\nu_1\ge\cdots\ge\nu_k=0$.
Let $\nabla$ be the $\mathcal O$-module of polynomials of column multidegree $\nu$ invariant under $X\mapsto XU$ for upper unitriangular $U$. Put
\[
 H=\prod_{j=1}^{k}\det(X_{[1,j],[1,j]})^{\nu_j-\nu_{j+1}},
\]
where $\nu_{k+1}=0$. Let $\Delta$ be the span of divided-power words applied to $H$.
Put $h=\nu_1$ and $V=K\nabla$.
Gaussian decomposition $X=LDU$ shows that the highest line is exactly $KH$: a highest vector is also invariant under lower unitriangular left multiplication.
Complete reducibility gives irreducibility of $V$; its weights have every coordinate between $0$ and $h$.
The coefficient lattice $\nabla$ is graded and divided-power stable, with highest line $\mathcal OH$ since $H(I)=1$.
The highest-weight projection $C$ therefore maps $\nabla$ into $\Delta$.
The Hom lattice $\operatorname{Hom}_{\mathcal O}(\nabla,\Delta)$ is divided-power stable, since its $j$th divided action is
$\sum_{a+b=j}(-1)^b E_{uv}^{[a]}T E_{uv}^{[b]}$.
Apply (IW.2) to $C$. Its invariant projection is $(\dim V)^{-1}1$; multiplication by the integer $\dim V$ proves
\[
 p^{D_p(k,h)}\nabla\subseteq\Delta\subseteq\nabla. \tag{IW.3}
\]
If $h=0$, both lattices equal $\mathcal O$. \textsf{[IW3]}

Define the \emph{symmetric bilinear} Fock form and its exact normalization by
\[
 F(P,Q)=\sum_\alpha P_\alpha Q_\alpha\alpha!,\qquad
 N=F(H,H)\in\mathbb Z_{>0},\qquad B=F/N.
\]
Write $R(g)P(X)=P(gX)$, the source's substitution convention.
Monomial calculation gives transpose adjointness; highest-weight projection and $H(I)=1$ give
\[
 B(R(g^t)H,Q)=B(H,R(g)Q)=Q(g)\qquad(Q\in V). \tag{IW.4}
\]
Thus $B(\Delta,\nabla)\subseteq\mathcal O$.
Conversely, if $B(\Delta,Q)\subseteq\mathcal O$, substitute a generic lower unitriangular $L$ in (IW.4).
The coefficients of $R(L^t)H$ are integral combinations of divided-power words, so $Q(L)$ is integral.
Gaussian decomposition and column homogeneity recover $Q(X)$ after inverting principal minors.
These minors are primitive polynomials; the DVR Gauss lemma removes their denominators.
Hence $Q\in\nabla$. Dualizing full lattices gives
\[
 \Delta^\#=\nabla,\qquad \nabla^\#=\Delta,
 \qquad p^{D_p(k,h)}\nabla\subseteq\nabla^\#\subseteq\nabla. \tag{IW.5}
\]
The expanded note proves the Gaussian recovery and nondegeneracy explicitly. \textsf{[IW5]}

\subsection{Tensor transport and the shared projector bound}

For a dominant integer weight $\lambda$, use $\nu_i=\lambda_i-\lambda_k$ and the determinant twist $\lambda_k$.
For a basis matrix $A$, the source convention is $L(A)=\det(A)^{-\lambda_k}R(A^{-1})\nabla$; for tuples use the tensor product of these lattices.
Transpose adjointness cancels the determinant twists at $A$ and $A^{-t}$.
Tensoring (IW.5), and identifying the linear dual with the opposite representation by the normalized Fock forms, proves
\[
 \begin{aligned}
 p^{\kappa}\,L(A^{-t})&\subseteq\Psi(L(A)^*)\subseteq L(A^{-t}),\\
 \kappa&=\sum_{s:\,h_s>0}\binom{k}{2}\lfloor\log_p(h_s+k-1)\rfloor.
 \end{aligned}
 \tag{IW.6}
\]
This is the actual scalar-normalized tensor pairing, not a choice of arbitrary perfect form. \textsf{[IW6]}
At an edge, the midpoint matrices are $A_s=\sigma^s(g)\operatorname{diag}(w^{-1}I_r,I)$ with $w^2=p$.
The opposite edge gives exactly $A_s^{-t}$; the common $w^{-A_r}$ factors in $R_e$ and $R'_e$ preserve (IW.6).

For completeness, (IW.2) also controls every isotypic projector $e_U$ on a graded, divided-power-stable full lattice $M$ in a split semisimple rational representation.
Here $h$ must bound $|\mu_a-\eta_a|$ for \emph{two} weights of that representation.
For an integral basis $u_i$ of $M\cap U$, let $\mathcal A$ be $I_h$ after weight-zero projection on $\operatorname{Hom}(U,V)$.
Then
\[
 \begin{aligned}
 e_U(y)&=(\dim U)\sum_i\mathcal A(u_i^*\otimes y)(u_i),\\
 p^{D_p(k,h)}e_U(M)&\subseteq M.
 \end{aligned}\tag{IW.7}
\]
Schur's trace identity proves the first formula on each irreducible summand; (IW.2) proves the second. \textsf{[IW4]}
Finally, $h_s+k\le a^2$ gives each exponent at most $k(k-1)\log_p a$.
For the two selected families this implies the source's stated $\kappa_0+\kappa_f\le4Nk^2\log_p a$. \textsf{[IW7]}
